\documentclass{article} % For LaTeX2e
\usepackage{iclr2025_conference,times}

% Optional math commands from https://github.com/goodfeli/dlbook_notation.
\input{math_commands.tex}

\usepackage{hyperref}
\usepackage{url}
\usepackage{booktabs}
\usepackage{xcolor}

\title{Optimizing Inter-node KV Cache Transfer in Prefill-Decode Disaggregation}

% Authors are now visible for the course project submission.
\author{Zhouxiang Fang, Yunfei Xie, Feng Luo, Haotian Xia \\
Department of Computer Science\\
Rice University\\
Houston, TX 77005, USA \\
\texttt{\{zf28, yx106, fl38, hx50\}@rice.edu} 
}

\newcommand{\fix}{\marginpar{FIX}}
\newcommand{\new}{\marginpar{NEW}}
\newcommand{\zf}[1]{{\color{blue} [{\bf Zhouxiang}: #1]}}
\newcommand{\hx}[1]{{\color{brown} [{\bf haotian}: #1]}}

\iclrfinalcopy % Uncommented for the final proposal version.
\begin{document}


\maketitle

\begin{abstract}
As Large Language Models (LLMs) scale, the traditional colocation of prefill and decode stages creates resource contention, resulting in suboptimal hardware utilization. 
While Prefill-Decode (PD) disaggregation decouples these phases to improve resource allocation, it introduces a significant bottleneck: the migration of the Key-Value (KV) cache. 
This project specifically targets the inter-node disaggregation regime, where network overhead becomes the dominant performance constraint. 
We plan to build an inter-node PD disaggregation framework and systematically evaluate three optimization strategies (KV cache compression, attention-based pruning, and chunked/pipelined transfer), both individually and in combination.
If time permits, we will also study the interaction effects among these techniques.
Our results will provide practical insights for deploying inter-node PD disaggregation systems.
\end{abstract}

\section{Introduction}

LLM serving pipelines exhibit a clear phase asymmetry. 
The prefill stage processes the prompt and produces the Key-Value (KV) cache; this stage is compute-heavy and its cost grows with context length. 
The decode stage generates tokens autoregressively and is dominated by KV cache reads and strict token-level latency requirements. 
In traditional colocated serving, these phases contend for the same GPU resources, and under bursty workloads, prefill spikes can stall active decode streams, inflating inter-token latency (ITL) and degrading tail latency.

\textbf{Prefill-Decode (PD) Disaggregation} addresses this by separating prefill and decode into distinct worker pools, enabling independent provisioning and scheduling~\citep{zhong2024distservedisaggregatingprefilldecoding, patel2024splitwiseefficientgenerativellm, qin2025mooncakekvcachecentricdisaggregatedarchitecture}. 
However, disaggregation requires migrating the KV cache from prefill workers to decode workers. 
The KV cache size grows linearly with prompt length and model depth: roughly $\text{KV bytes} \approx 2 \times L \times S \times H \times D \times b$, where $L$ is the number of layers, $S$ the sequence length, $H$ the number of attention heads, $D$ the head dimension, and $b$ the bytes per element (e.g., $b{=}2$ for FP16).
For long-context prompts, this payload can reach gigabytes per request.
This migration cost is minor within a single node, but becomes the dominant bottleneck in the \textbf{inter-node} regime, where performance becomes sensitive to network bandwidth, latency, and congestion.

In this project, we focus on \textbf{inter-node KV transfer optimization} for PD disaggregation. 
Existing approaches reduce transfer overhead via (i) KV compression/quantization~\citep{liu2024cachegenkvcachecompression, zhang2025hackhomomorphicaccelerationcompression}, (ii) attention-based pruning/eviction~\citep{zhang2023h2oheavyhitteroracleefficient, zhang2025pdtrimtargetedpruningprefilldecode}, or (iii) chunked/pipelined transfer engines~\citep{qin2025mooncakekvcachecentricdisaggregatedarchitecture, kwon2023efficient}.
Yet these are largely studied in isolation.
In practice, a serving system may combine them, but the interaction effects (e.g., pruning changing quantization error profiles, quantization affecting pipeline granularity and amortization) remain under-characterized. 
Our goal is to systematically evaluate these strategies and their combinations under controlled network conditions.

\section{Methodology}

\subsection{Framework and Proposed Experiments}

We plan to build an inter-node PD prototype on top of vLLM~\citep{kwon2023efficient}. 
Requests are processed by a \emph{prefill GPU} that computes KV blocks and then forwarded to a \emph{decode GPU} that performs autoregressive generation. 
The KV transfer path will use RDMA (InfiniBand or RoCE) when available on the cluster. For reproducibility, we will also vary effective bandwidth and latency using standard traffic-control tools (e.g., Linux \texttt{tc}/\texttt{netem}) over TCP/IP to test non-RDMA regimes.

On top of a vanilla inter-node PD baseline (raw KV transfer), we will add three optimizations with configurable parameters:
\begin{itemize}
    \item Compression/quantization (e.g., FP16/BF16 baseline vs.\ 8-bit/4-bit variants with measured (de)quantization cost)
    \item Attention-based pruning/eviction (remove a tunable fraction of low-importance tokens before transfer, e.g., via attention-score-based selection)
    \item Chunked/pipelined transfer (stream KV blocks in chunks to overlap network transfer with early decode computation).
\end{itemize}
We will evaluate each technique in isolation and in combination under matched settings to quantify interaction effects.

Workloads will include prompt-heavy (long input/short output), generation-heavy (short input/long output), and long-context chat-like mixes. We will also vary the concurrency level.

\subsection{Evaluation Metrics}

We will report: \textbf{TTFT}, mean and tail \textbf{ITL} (P95/P99), and request-level latency/throughput under load. We will quantify \textbf{SLO attainment} under a latency budget and compute \textbf{break-even} curves over bandwidth/RTT where inter-node PD becomes worse than colocated serving. We will also instrument an \textbf{overhead breakdown} (compute, (de)compression, pruning, transfer, and achieved overlap). Finally, we will measure \textbf{quality impact} of pruning/quantization using perplexity and/or accuracy on small representative tasks.

\subsection{Implementation Details and plans}
We will instrument the PD pipeline into the following stages:
(1) prefill compute,
(2) (optional) pruning decision and KV compaction,
(3) (optional) KV (de)compression/(de)quantization,
(4) KV serialization and transport (RDMA when available; otherwise TCP with controlled netem),
and (5) decode warm-up and streaming generation.
For pipelining, we will support configurable chunk sizes and measure achieved overlap between communication and prefill computation.
We will report an overhead breakdown for each stage to explain latency changes beyond end-to-end metrics.

\section{Expected Results and Contributions}

We expect the benefit of inter-node PD disaggregation to depend on the network conditions and workload characteristics. In prompt-heavy regimes, TTFT should be most sensitive to transfer delay; in generation-heavy regimes, ITL tail behavior should dominate. We anticipate that pipelining becomes necessary below certain bandwidth thresholds, and that combining it with compression and pruning may produce non-obvious trade-offs due to chunk granularity and amortization.

Our main contribution will be a \textbf{systematic evaluation} of compression, pruning, and pipelining in a unified inter-node PD pipeline. Deliverables include: (i) a reproducible benchmark harness with controlled network conditions, (ii) measurements showing how each optimization (and their combinations) affect TTFT, ITL tail latency, and SLO attainment over vanilla inter-node PD, and (iii) if time permits, an analysis of interaction effects among the three strategies.



\bibliography{iclr2025_conference}
\bibliographystyle{iclr2025_conference}

\section*{Appendix: AI Usage Disclosure}

In accordance with the course requirements, we disclose the use of AI tools (Gemini 3 Flash) in the development and refinement of this project proposal. AI was utilized as a collaborative partner in the following capacities:

\begin{itemize}
    \item \textbf{Literature Review}: AI assisted in summarizing and contrasting current Prefill-Decode (PD) disaggregation frameworks (e.g., DistServe, Splitwise). This helped us pinpoint inter-node KV cache migration as the primary performance constraint when scaling to distributed clusters.
    \item \textbf{Methodology Refinement}: AI provided technical suggestions for simulating diverse network environments using Linux traffic-control tools (\texttt{tc/netem}) to ensure our results remain reproducible in non-RDMA regimes.
    \item \textbf{Metrics Selection}: The conceptual framework for calculating "break-even curves" (identifying the network thresholds where disaggregation outperforms colocation) was developed through iterative brainstorming sessions with the AI tool.
\end{itemize}

\end{document}